\section{Podsumowanie i wnioski}
\label{sec:podsumowanie}

\subsection{Realizacja celu}

Zrealizowano wszystkie wymagania funkcjonalne z rozdziału~\ref{sec:cel}.
Kierowca zakłada konto i loguje się, a żądania są uwierzytelniane tokenem JWT.
Uzupełnia profil, do którego przypisane są jego dokumenty. Dokument fotografuje
aparatem telefonu, kadruje i obraca zdjęcie, a aplikacja kompresuje je przed
wysłaniem (zdjęcie 12~Mpx o rozmiarze 3,7~MB zajmuje po kompresji około
0,7--0,8~MB). Dokumenty mają typy, statusy z regułami przejść, kolejne wersje
pliku i historię zmian. Można je wyszukiwać i filtrować.

Wymagania niefunkcjonalne zrealizowano następująco:
\begin{itemize}
  \item \textbf{Bezpieczeństwo}: pliki są szyfrowane algorytmem AES-256-GCM
        z osobnym kluczem dla każdego pliku (szyfrowanie kopertowe), hasła są
        przechowywane jako skróty argon2id, a nieudane logowania ograniczane.
        Użytkownik ma dostęp tylko do własnych dokumentów, co sprawdzają testy.
  \item \textbf{Obsługa mobilna}: responsywny interfejs z dolną nawigacją,
        instalowalny jako PWA i uruchamiający się bez sieci.
  \item \textbf{Łatwe wdrożenie}: cały system, łącznie z monitoringiem,
        uruchamia polecenie \texttt{docker compose up}. Migracje bazy
        wykonuje osobny kontener przed startem backendu.
  \item \textbf{Obserwowalność i niezawodność}: metryki Prometheus, cztery
        dashboardy Grafany, 12 reguł alertów, centralne logi w Loki oraz
        automatyczne restarty po awarii i po zawieszeniu procesu, sprawdzone
        testem chaosu.
  \item \textbf{Jakość}: 158 testów backendu (głównie funkcjonalnych, na
        prawdziwej bazie i magazynie plików), 90 testów pakietu wspólnego,
        46 testów jednostkowych frontendu i 31 scenariuszy Playwright
        w emulacji telefonu. Pokrycie kodu przekracza 93~\% instrukcji
        w każdym pakiecie (frontend mierzony testami w przeglądarce).
\end{itemize}

\subsection{Napotkane problemy i wnioski}

Najcenniejszym wnioskiem z projektu jest skuteczność testów funkcjonalnych.
Wykryły one kilka błędów, których nie ujawniłyby testy jednostkowe ani
pobieżne ręczne sprawdzenie:
\begin{itemize}
  \item po odświeżeniu strony pierwsze zapytania wysyłane były bez tokenu,
        bo efekty komponentów potomnych w React wykonują się przed efektem
        dostawcy kontekstu (poprawka: token przechowywany poza drzewem
        komponentów);
  \item strumień aparatu wyłączał się zaraz po uruchomieniu, a użytkownik
        widziałby zamrożoną klatkę (poprawka: efekt uruchamiany raz przy
        wyświetleniu widoku);
  \item po wylogowaniu i zalogowaniu innego kierowcy aplikacja wracała na
        stronę dokumentu poprzedniego użytkownika;
  \item metryki HTTP miały błędne etykiety tras, bo Express przywraca
        ścieżkę bazową routera przed zakończeniem odpowiedzi;
  \item pomiar opóźnienia pętli zdarzeń tracił próbkę po wyzerowaniu
        histogramu, więc zawieszenie procesu nie byłoby wykryte.
\end{itemize}
Każdy z tych błędów najpierw odtworzono testem, który nie przechodził,
a dopiero potem poprawiono kod. Testy pozostają w zestawie i chronią przed
powrotem błędu.

Drugi wniosek dotyczy środowiska testowego. Testy w przeglądarce okazały się
wrażliwe na obciążenie komputera: haszowanie haseł argon2 i przetwarzanie
zdjęć 12~Mpx konkurowały o procesor z przeglądarkami. Stabilność zapewniło
uruchamianie testów na zbudowanej aplikacji (zamiast serwera deweloperskiego)
i kolejne, a nie równoległe, wykonywanie scenariuszy.

Trzeci wniosek dotyczy monitoringu: sama konfiguracja kontroli stanu nie
gwarantuje samonaprawy. Test chaosu pokazał, że Docker nie restartuje
kontenera zatrzymanego poleceniem \texttt{docker kill}, bo traktuje je jak
ręczne zatrzymanie, a zawieszonego, ale działającego procesu nie restartuje
wcale. Dopiero dodatkowy kontener \emph{autoheal} i test z rzeczywistym
zawieszeniem procesu potwierdziły, że system sam wraca do działania.

\subsection{Ograniczenia}

Aparat sprawdzono tylko w emulacji (wirtualna kamera Chromium), bez
prawdziwego telefonu. Procedura testu na urządzeniu z Androidem jest opisana
w pliku README. Potok ciągłej integracji jest przygotowany, ale nie był
uruchomiony, bo repozytorium nie ma zdalnej kopii. Aplikacja ma jedną rolę
użytkownika, więc zmiany statusu, które w praktyce należałyby do biura
(akceptacja, odrzucenie), wykonuje sam kierowca.

\subsection{Możliwości rozwoju}

\begin{itemize}
  \item \textbf{Lepsze zdjęcia dokumentów}: automatyczne wykrywanie krawędzi
        kartki i korekcja perspektywy, jak w aplikacjach skanujących
        (podrozdział~\ref{sec:przeglad}).
  \item \textbf{Rozpoznawanie tekstu (OCR)}: odczyt numeru dokumentu,
        nadawcy i odbiorcy ze zdjęcia i wypełnianie nimi formularza.
  \item \textbf{Panel dyspozytora}: rola pracownika biura, który akceptuje
        lub odrzuca dokumenty kierowców i widzi dokumenty całej floty.
  \item \textbf{e-CMR i eFTI}: import i eksport elektronicznych listów
        przewozowych oraz integracja z platformami eFTI, które od 2027~r.
        mają być w pełni honorowane przez organy kontrolne.
  \item \textbf{Praca bez sieci}: kolejka wysyłek zapisywana na urządzeniu
        w postaci zaszyfrowanej i wysyłana po odzyskaniu połączenia, co ma
        znaczenie na trasach bez zasięgu.
  \item \textbf{Zarządzanie kluczami}: przechowywanie klucza głównego
        w usłudze KMS i jego okresowa rotacja (szyfrowanie kopertowe
        pozwala zmienić klucz główny bez ponownego szyfrowania plików).
\end{itemize}
